
% rappel liste
% rappel dictionnaire
% rappel fonctions

%tyoe variable, parcours, tour à tour, résultat de la variable
\setcounter{chapter}{-1}
\chapter*{Rappel de 1er NSI}
 Nous allons étudier l'implémentation du morpion. Le jeu est implémenté selon le déroulement suivant : Chaque joueur joue chacun son tour. Le programme demande au joueur où il souhaite jouer. La grille est mise à jour à chaque action. Le jeu s’arrête si la grille est pleine ou si quelqu’un a gagné.

Voici la fonction qui affiche la grille du morpion :
\begin{lstlisting}{style=pythonstyle, language=Python}
# La fonction qui affiche la grille dans le terminal
def afficherGrille(g, gagnant=None):
    print(" " + g[0][0] + " | " + g[0][1] + " | " + g[0][2] + " ")
    print("---+---+---")
    print(" " + g[1][0] + " | " + g[1][1] + " | " + g[1][2] + " ")
    print("---+---+---")
    print(" " + g[2][0] + " | " + g[2][1] + " | " + g[2][2] + " \n")
    
    if gagnant:
        print("* Partie terminée : le joueur", gagnant, "a gagné. *")


grille = [[" " for _ in range(3)] for _ in range(3)]  # crée un tableau de 9 caractères espaces " 
\end{lstlisting}

\begin{Exercise}[title={Rappels de fonction}]
\Question Quel mot clef en python permet de reconnaître une définition de fonction ?
\Question Quels sont les paramètres de la fonction afficherGrille ?
\Question Dans le code qui vous a été fourni, quel est le type de grille ? À quoi le reconnaissez-vous ?
\Question Écrivez l'appel de fonction de afficherGrille avec uniquement la variable \textit{grille} en paramètre.
\Question Qu'affiche sur le terminal l'appel suivant afficherGrille(grille, gagnant="X")
\end{Exercise}

\begin{lstlisting}{style=pythonstyle, language=Python}
joueurs = ("X","O")
\end{lstlisting}
\begin{Exercise}[title={Rappels sur les tuples}]
\Question Quel est le type de la variable joueurs ?
\Question A votre avis que désigne-t-elle ?
\Question Écrire la ligne de code qui affiche le nom du deuxième joueur.
\Question Pouvez-vous modifier le nom du deuxième joueur ? Justifier.
\end{Exercise}


\begin{lstlisting}{style=pythonstyle, language=Python}
def tourJoueur(joueur, grille):
    tour = _ _ _ _ input("Sur quelle case voulez vous jouer ?")
    grille[_ _ _ _ _ ] = joueur
\end{lstlisting}

\vfill
% \begin{lstlisting}{style=pythonstyle, language=Python}
% def grillePleine(grille):
%     for elt in grille :
%         if elt == " ":
%             return False
%     return True
% \end{lstlisting}

% \begin{Exercise}[title={Rappels boucle for}]
% \Question Est-ce un parcours par index ou par valeur ?
% \Question Quel est le type des retours ? Que représentent-ils ?
% \end{Exercise}


\begin{Exercise}[title={Rappels sur les listes}]
\begin{lstlisting}{style=pythonstyle, language=Python}
def ligneGagnante(grille, ligne):
    if grille[ligne][0] == grille[ligne][1] == grille[ligne][2]:
        return grille[ligne][0]
    else : return False
    
def colonneGagnante(grille, colonne):
    if grille[0][colonne] == ___________________ == ___________________:
        return grille[0][colonne]
    else : return False

def diagonalGagnante(grille):
    if ___________________ == ___________________== ___________________ :
        return grille[0][0]
    elif ___________________ == ___________________ == ___________________:
        return grille[2][0]
    else : return False

\end{lstlisting}
%len ?
\Question Complétez les trous
\Question Quelles sont les valeurs que peut retourner cette fonction ?

\end{Exercise}

\vfill

\begin{Exercise}[title={Rappels sur les dictionnaires}]
\begin{lstlisting}{style=pythonstyle, language=Python}
joueurs = {"Lagon": (4, 15),
           "Corail": (6, 15),
           "Manta": (11, 20),
           "Tortue": (19, 30)}

def meilleurRatio(joueurs):
    meilleur = None
    rmax = 0
    for n, (v, j) in joueurs.items():
        if v/j > rmax:
            rmax = v/j
            meilleur = n
    return meilleur            
\end{lstlisting}

Le dictionnaire joueurs a pour clé le nom du joueur et pour valeur un tuple (nombre de partie gagnée, nombre de parties jouées).\\
On pourra utiliser la fonction meilleurRatio pour trouver le joueur ayant le meilleur ratio de victoire.\\
Pour trouver le total de parties jouées, on peut additionner les deuxièmes valeurs des tuples.
\Question Quelles sont les clés du dictionnaire joueurs ?
\Question Quelle est le type des valeurs du dictionnaire joueurs ?
\Question Que vaut joueurs["Tortue"][0], a quoi cela correspond t-il ?
\Question A quoi correspond joueurs["Lagon"][1] + joueurs["Corail"][1] + joueurs["Manta"][1] \\ + joueurs["Tortue"][1] ?
\Question A quoi correspondent les variables n, v, j de la fonction meilleurRatio ?
\Question Quel est le résultat de meilleurRatio(joueurs) ?
\Question En cas d'égalité sur le ratio, quel joueur sera retourné par la fonction.
\Question Tortue vient de gagner une partie contre Lagon donner les instructions permettant de mettre à jour le dictionnaire joueurs.
\end{Exercise}

\newpage
%Exercice sur les listes : 1) affichez des valeurs, 2) parcourir des listes, 3) Ajouter les valeurs à une liste

% \begin{Exercise}
%     \begin{lstlisting}
% L = [8, 24, 48, 2, 16]
% taille = len(L)

% #nombre de multiples de 3 présents dans L
% mult3 = 0
% for i in range(_ _ _ _ _ ) :
%  if _ _ _ _ _ _ _ _ _  == 0 :
%       mult3 += _ _ _
% print ("Nombre de multiples de 3 : ", mult3)
%     \end{lstlisting}
% \Question Complétez le code à trou ci-dessus.
% \end{Exercise}


\begin{Exercise}[title = {Listes et dictionnaires}]
Voici une structure de données qui récupère le nombres de médailles gagnées par 10 pays participant aux JO de Paris 2024.
\begin{lstlisting}
JO_2024 =[["pays" , "or", "argent", "bronze" ] ,
        ["Italie", 12, 13, 15],
        ["Japon", 20, 12, 13],
        ["Etats-Unis", 40, 44, 42],
        ["Pays-Bas", 15, 7, 12],
        ["France", 16, 26, 22],
        ["Chine",40, 27, 24],
        ["Corée du Sud", 13, 9, 10],
        ["Allemagne", 12, 13, 8],
        ["Australie", 18, 19, 16],
        ["Grande-Bretagne", 14, 22, 29]]
\end{lstlisting}
\Question Quelle est la structure de données choisie ?
\Question Donner l'instruction pour afficher le premier pays de la liste.
\Question Donner l'instruction pour afficher le dernier pays de la liste.
\Question Donner l'instruction pour connaître le nombre de pays contenus dans la liste.
\Question Le gagnant est celui qui a gagné le plus de médailles d'or. Écrire une fonction \textit{gagnant} qui prend en paramètre la liste et qui renvoie le nom du pays qui a gagné le plus de médailles d'or ainsi que le nombre de médailles d'or.\textit{Vous allez devoir parcourir la liste et trouver le maximum de médailles d'or}
\Question Écrire une fonction \textit{plus\_medailles} qui prend en paramètre la liste et qui renvoie le nom du pays qui a gagné le plus de médailles ainsi que le nombre total de médailles gagnées.

\Question Proposez une autre structure de données pour l'enregistrement du nombre de médailles gagnées par les pays participant aux JO.
\Question Ecrivez votre nouvel enregistrement avec votre nouvelle structure de données.
\Question Répondez aux mêmes questions que l'exercice précédent.
\Question A votre avis, quelle structure est la plus adaptée ?
\end{Exercise}
